\section{Mode Memori Berpenunjuk}
\label{sec:mode-berpenunjuk}

Empat mode berikutnya menyatakan alamat melalui register penunjuk,
bukan bilangan tetap. Karena isi register dapat diubah, alamatnya
dapat berubah saat program berjalan. Kemampuan inilah yang membuat
akses berulang terhadap array dan \eng{record} menjadi mungkin.

\subsection{Indirect Memory}

Pada mode \term{indirect memory}, operand ditulis di antara tanda
kurung siku yang berisi satu register. Register itu tidak menyimpan
data, melainkan \eng{offset} tempat data berada.

\begin{lstlisting}[language=MASM, numbers=none]
MOV BX, OFFSET data1    ; BX <- offset awal array
MOV AX, [BX]            ; AX <- isi memori DS:BX
ADD BX, 2               ; maju dua byte (satu elemen word)
MOV AX, [BX]            ; AX <- elemen kedua
\end{lstlisting}

Hanya \reg{BX}, \reg{BP}, \reg{SI}, dan \reg{DI} yang boleh menjadi
penunjuk alamat. \reg{BX} dan \reg{BP} disebut \term{register basis},
sedangkan \reg{SI} dan \reg{DI} disebut \term{register indeks}.

\subsection{Base Relative dan Direct Index}

Bentuk \code{[BX+disp]} dan \code{[BP+disp]} disebut mode
\term{base relative}. Bentuk \code{[SI+disp]} dan \code{[DI+disp]}
disebut mode \term{direct index}.

\begin{lstlisting}[language=MASM, numbers=none]
; base relative
MOV BX, OFFSET data1
MOV AX, [BX+4]          ; elemen ke-2 dari array word

; direct index
MOV SI, 4
MOV AX, data1[SI]       ; DS:data1 + SI
MOV AX, [SI+8]          ; DS:SI + 8
\end{lstlisting}

Penulisan \code{data1[SI]} diterjemahkan assembler menjadi
\code{[data1 + SI]}, menyerupai penulisan array pada bahasa tingkat
tinggi. Perpindahan boleh berukuran 8 bit, 16 bit, atau nol.

\subsection{Based Indexed}

Mode \term{based indexed} menggabungkan satu basis, satu indeks, dan
satu perpindahan sekaligus. Hanya empat kombinasi yang legal:
\code{[BX+SI+disp]}, \code{[BX+DI+disp]}, \code{[BP+SI+disp]}, dan
\code{[BP+DI+disp]}.

\begin{lstlisting}[language=MASM, numbers=none]
MOV BX, OFFSET tabel    ; basis: awal tabel record
MOV SI, 6               ; indeks: record ke-3 (3 x 2 byte)
MOV AX, [BX+SI+2]       ; medan ke-2 dari record ke-3
\end{lstlisting}

Mode ini dipakai bila alamat akhir tersusun dari dua komponen yang
bergerak sendiri-sendiri, misalnya baris dan kolom matriks, atau
nomor \eng{record} dan nomor medannya.

\subsection{Alamat Efektif}

Prosesor membentuk \term{alamat efektif} (\eng{effective address},
EA), yaitu jumlah basis, indeks, dan perpindahan. Hasilnya berupa
offset 16 bit. Alamat fisik dibentuk dengan menambahkan basis segmen.

\begin{figure}[H]
    \centering
    \begin{tikzpicture}[
        komp/.style={
            draw=obeBlue, line width=0.8pt, fill=obeBlueLight,
            rounded corners=3pt, minimum width=3.9cm,
            minimum height=1.1cm, align=center, inner sep=4pt,
            font=\small
        },
        jumlah/.style={
            draw=obeGray, line width=0.9pt, fill=obeGrayLight,
            circle, minimum size=1.0cm, font=\large\bfseries
        },
        hasil/.style={
            draw=obeTeal, line width=1.0pt, fill=obeTealLight,
            rounded corners=3pt, minimum width=4.2cm,
            minimum height=1.1cm, align=center, inner sep=4pt,
            font=\small
        },
        panah/.style={-{Stealth[length=2.5mm]}, line width=0.9pt,
            obeBlue},
        kecil/.style={font=\scriptsize\itshape, text=obeGray}
    ]
        \node[komp] (base) at (0, 2.5)
            {\textbf{Basis}\\{\scriptsize \reg{BX}, \reg{BP}, atau
             tidak ada}};
        \node[komp] (index) at (0, 1.2)
            {\textbf{Indeks}\\{\scriptsize \reg{SI}, \reg{DI}, atau
             tidak ada}};
        \node[komp] (disp) at (0, -0.1)
            {\textbf{Perpindahan}\\{\scriptsize 8 bit, 16 bit, atau
             0}};

        \node[jumlah] (plus) at (4.7, 1.2) {$+$};

        \node[hasil] (ea) at (8.6, 1.2)
            {\textbf{Alamat efektif}\\[1pt]
             {\scriptsize \code{EA = basis + indeks + disp}}};

        \draw[panah] (base) -- (plus);
        \draw[panah] (index) -- (plus);
        \draw[panah] (disp) -- (plus);
        \draw[panah] (plus) -- (ea);

        \node[hasil, below=1.3cm of ea] (fisik)
            {\textbf{Alamat fisik}\\[1pt]
             {\scriptsize \code{(segmen x 16) + EA}}};
        \draw[panah] (ea) -- (fisik)
            node[midway, right=3pt, kecil] {ditambah basis segmen};
    \end{tikzpicture}
    \caption{Perhitungan alamat efektif dan alamat fisik.}
    \label{fig:alamat-efektif}
\end{figure}

Misalkan \reg{DS} = \hex{1000}, \reg{BX} = \hex{0200},
\reg{SI} = \hex{0006}, dan instruksi
\lstasm|MOV AX, [BX+SI+2]|. Alamat efektifnya
$\hex{0200} + \hex{0006} + \hex{0002} = \hex{0208}$, dan alamat
fisiknya $\hex{1000} \times \hex{10} + \hex{0208} = \hex{10208}$.

\subsection{Kombinasi Register yang Legal}

Tidak semua penulisan operand berpenunjuk dapat diterima assembler.
Tabel~\ref{tab:kombinasi-register} merangkum kombinasinya.

\begin{table}[H]
    \centering
    \caption{Kombinasi register penunjuk alamat pada 8086.}
    \label{tab:kombinasi-register}
    \small
    \begin{tabular}{@{}p{3.4cm}p{1.6cm}p{1.6cm}p{6.8cm}@{}}
        \toprule
        \textbf{Bentuk operand} & \textbf{Legal} &
        \textbf{Segmen} & \textbf{Alasan} \\
        \midrule
        \code{[BX]}, \code{[SI]}, \code{[DI]} & Ya & \reg{DS} &
        Tersedia untuk satu penunjuk \\
        \addlinespace
        \code{[BP]} & Ya & \reg{SS} &
        \reg{BP} adalah basis segmen \eng{stack} \\
        \addlinespace
        \code{[BX+SI]}, \code{[BX+DI]} & Ya & \reg{DS} &
        Satu basis ditambah satu indeks \\
        \addlinespace
        \code{[BP+SI]}, \code{[BP+DI]} & Ya & \reg{SS} &
        Basis \reg{BP} ditambah satu indeks \\
        \addlinespace
        \code{[AX]}, \code{[CX]}, \code{[DX]} & Tidak & --- &
        Bukan basis maupun indeks \\
        \addlinespace
        \code{[SP]} & Tidak & --- &
        \reg{SP} bukan basis; gunakan \reg{BP} \\
        \addlinespace
        \code{[BX+BP]} & Tidak & --- &
        Dua register basis tidak dapat digabung \\
        \addlinespace
        \code{[SI+DI]} & Tidak & --- &
        Dua register indeks tidak dapat digabung \\
        \addlinespace
        \code{[BX+SI+DI]} & Tidak & --- &
        Maksimal satu basis ditambah satu indeks \\
        \bottomrule
    \end{tabular}
\end{table}

\begin{peringatanbox}
\vspace{1.1em}
Penulisan seperti \code{[AX]} ditolak assembler karena 8086 tidak
menyediakan cara pengalamatan dengan \reg{AX} sebagai penunjuk.
Gunakan \reg{BX}, \reg{BP}, \reg{SI}, atau \reg{DI}.
\end{peringatanbox}
